\subsection{Признаки}

Центральным элементом данной работы является не выбор конкретной модели, а сравнение нескольких групп признаков, используемых для прогноза направления движения цены. В машинном обучении качество прогноза зависит не только от алгоритма, но и от того, какую информацию содержат входные переменные. Поэтому признаки рассматриваются как основной объект эмпирической проверки: необходимо определить, дают ли текстовые признаки, извлечённые с помощью FinBERT, самостоятельный или добавочный сигнал по сравнению с рыночными признаками \cite{guyon2003feature}.

В эксперименте используются три группы признаков: \texttt{prices-only}, \texttt{FinBERT-only} и \texttt{prices + FinBERT}. Такое разделение позволяет провести ablation-сравнение, то есть оценить вклад каждой группы признаков отдельно и в комбинации. В литературе по финансовому NLP похожая логика используется в мультимодальных постановках, где сравниваются ценовые, текстовые и объединённые источники информации \cite{hiew2019bert,zou2022prebit,albada2025bertweet}.

\begin{figure}[htbp]
\centering
\begin{tikzpicture}[
    node distance=1.2cm,
    box/.style={draw, rounded corners, align=center, minimum width=3.2cm, minimum height=0.8cm},
    arrow/.style={-{Latex[length=2mm]}, thick}
]
\node[box] (prices) {Рыночные данные\\OHLCV, returns};
\node[box, right=1.4cm of prices] (text) {Финансовые тексты\\news, posts};
\node[box, below=1.0cm of prices] (pricefeat) {prices-only\\$p_t \in \mathbb{R}^{11}$};
\node[box, below=1.0cm of text] (textfeat) {FinBERT-only\\$e_t \in \mathbb{R}^{768}$};
\node[box, below right=1.1cm and -1.6cm of pricefeat] (combined) {combined\\$x_t=[p_t;e_t]$};
\node[box, below=1.0cm of combined] (model) {модель классификации\\$\hat y_t=P(y_t=1\mid x_t)$};
\draw[arrow] (prices) -- (pricefeat);
\draw[arrow] (text) -- (textfeat);
\draw[arrow] (pricefeat) -- (combined);
\draw[arrow] (textfeat) -- (combined);
\draw[arrow] (combined) -- (model);
\end{tikzpicture}
\caption{Общая схема формирования и сравнения групп признаков}
\end{figure}

\subsubsection{Рыночные признаки}

Первая группа признаков --- рыночные признаки. Они строятся на основе табличных данных, отражающих динамику цены актива. В работе используются следующие переменные:

\[
open,\ high,\ low,\ close,\ adj\text{-}close,\ inc\text{-}5,\ inc\text{-}10,\ inc\text{-}15,\ inc\text{-}20,\ inc\text{-}25,\ inc\text{-}30.
\]

Формально вектор рыночных признаков для дня $t$ обозначим как

\[
p_t = (open_t, high_t, low_t, close_t, adjclose_t, inc5_t, inc10_t, inc15_t, inc20_t, inc25_t, inc30_t) \in \mathbb{R}^{11}.
\]

Признаки \texttt{open}, \texttt{high}, \texttt{low}, \texttt{close} и \texttt{adj-close} описывают состояние цены внутри торгового дня и цену закрытия с учётом корректировок. Признаки \texttt{inc-5}, \texttt{inc-10}, \texttt{inc-15}, \texttt{inc-20}, \texttt{inc-25}, \texttt{inc-30} отражают относительные изменения за несколько предшествующих периодов. В совокупности эта группа признаков описывает краткосрочную ценовую динамику и используется как базовая рыночная информация.

Теоретически рыночные признаки могут содержать сигнал, связанный с инерцией, краткосрочными разворотами, волатильностью и реакцией цены на предшествующие события. Однако они описывают только уже произошедшее движение цены. Если часть информации о будущей динамике содержится в текстовом фоне, то одной ценовой истории может быть недостаточно.

\subsubsection{Текстовые признаки FinBERT}

Вторая группа признаков --- текстовые признаки. Они получены из финансовых текстов с помощью модели FinBERT. В данной работе FinBERT используется как генератор embedding-признаков. Каждый текст преобразуется в числовой вектор размерности $768$:

\[
e_t = FinBERT(text_t) \in \mathbb{R}^{768}.
\]

Такой вектор представляет текст не как набор отдельных слов, а как контекстное представление, в котором учитываются смысловые связи внутри сообщения. В отличие от простого sentiment-score, embedding может содержать более широкую информацию: упоминания компаний, событий, рыночных ожиданий, неопределённости, характера новости и общего контекста. Это соответствует переходу от словарных подходов к контекстным языковым моделям, описанному в работах по BERT и FinBERT \cite{devlin2019bert,araci2019finbert,yang2020finbert}.

Теоретическое основание использования FinBERT-признаков состоит в том, что финансовые тексты могут отражать информацию, ещё не полностью выраженную в ценовой динамике. Если такая информация систематически связана с последующим движением актива, embedding-признаки должны улучшать прогноз относительно модели без текстов. При этом литература показывает, что выигрыш от текстовых признаков обычно зависит от источника текста, горизонта прогноза и качества временного выравнивания \cite{tetlock2007media,loughran2011liability,hiew2019bert,zou2022prebit}.

\subsubsection{Объединённые признаки}

Третья группа --- объединённые признаки. Она включает как рыночные показатели, так и embedding-представления текстов. Формально объединённый вектор можно записать как

\[
x_t = [p_t; e_t] \in \mathbb{R}^{779},
\]

где $p_t \in \mathbb{R}^{11}$ --- рыночные признаки, а $e_t \in \mathbb{R}^{768}$ --- FinBERT embedding. Эта группа является ключевой для проверки основной гипотезы работы: дают ли текстовые признаки дополнительную информацию поверх уже имеющейся ценовой динамики.

Если FinBERT-признаки действительно содержат новый сигнал, то модель на \texttt{prices + FinBERT} должна показывать качество выше, чем модель на \texttt{prices-only}. В этом случае можно говорить о добавочной предсказательной силе текстовой информации. Если же объединённая модель не улучшает качество относительно \texttt{prices-only}, то это означает, что либо текстовые признаки не содержат полезного сигнала для выбранной задачи, либо этот сигнал не извлекается текущим способом, либо он уже косвенно отражён в рыночных признаках.

\subsubsection{Временное выравнивание признаков}

Для финансовых задач особенно важно, чтобы признаки были доступны до момента прогноза. Если модель использует текст или цену, появившиеся после прогнозируемого события, возникает утечка данных. В данной работе используется консервативное правило: для прогноза после дня $t$ текстовые признаки берутся не позднее $t-1$.

Целевая переменная определяется как:

\[
y_t = \mathbf{1}\{Close_{t+1} > Close_t\}.
\]

Следовательно, признаки должны описывать информацию, доступную до $Close_{t+1}$. Такой подход позволяет интерпретировать эксперимент как проверку прогностической, а не ретроспективной связи.

\subsubsection{Логика проверки влияния признаков}

Влияние признаков оценивается через сравнение качества моделей на разных feature groups. Для формализации добавочной силы текстовых признаков можно рассматривать разность метрик:

\[
\Delta_{text} = M(\texttt{prices + FinBERT}) - M(\texttt{prices-only}),
\]

где $M$ --- выбранная метрика качества. Если $\Delta_{text} > 0$, то текстовые признаки улучшают результат относительно рыночного baseline по данной метрике. Если $\Delta_{text} \leq 0$, то добавочная предсказательная сила текстовой группы в данном протоколе не подтверждается.

\begin{table}[htbp]
\centering
\begin{tabular}{ll}
\toprule
Сравнение & Что проверяется \\
\midrule
\texttt{FinBERT-only} vs \texttt{majority baseline} & самостоятельный текстовый сигнал \\
\texttt{prices-only} vs \texttt{majority baseline} & сигнал в рыночных признаках \\
\texttt{prices + FinBERT} vs \texttt{prices-only} & добавочная сила текстовых признаков \\
\texttt{prices + FinBERT} vs \texttt{FinBERT-only} & эффект объединения признаков \\
LSTM vs Logistic Regression & зависимость от последовательной архитектуры \\
\bottomrule
\end{tabular}
\caption{Логика сравнения групп признаков}
\end{table}

Таким образом, экспериментальная часть работы построена не как поиск лучшей модели, а как проверка информативности разных групп признаков. Главный исследовательский вопрос состоит в том, улучшают ли FinBERT-generated features прогноз направления цены и проявляется ли их добавочная сила при объединении с рыночными признаками.
